\section{限制与结论}

本文存在以下局限。全部验证基于单台标定系统上的单颗芯片（248 帧），缺乏同场景的高分辨率热成像参照；所有质量判断均来自内部一致性门控，而非外部真值比对。代理指标能够检出异常漂移，但不能直接认证恢复细节；\eng{PSF} 标定的不确定度进一步将可信交付限定在轮廓可见性层面。诊断框架向多系统、多场景的推广是自然的后续方向。

在整个采集与评估链中都不存在高分辨率真值的工业 \eng{LWIR} 微扫描场景下，本文揭示了学习类多帧超分辨的一种系统性失效：网络新增的高频细节沿观测算子零空间漂移，前向一致性损失对此近乎不可见。这意味着"看起来更锐利"与"重建更准确"之间存在不可由观测侧损失弥合的系统性脱节。实验表明，经典变分方法与学习类方法呈现互补的失效模式——前者观测保真但受限于正则化假设，后者表观丰富但高频缺乏观测支撑；在无外部真值的条件下，不存在可被单独认证的最优方法。
